\section{Related Work}
\label{sec:related_work}

The connection between inconsistency and planning has been explored from several perspectives, with foundational work establishing important theoretical and diagnostic approaches for unsolvable planning problems. While the broader field of inconsistency measurement in propositional logic has developed various techniques for quantifying conflicts in logical knowledge bases~\cite{Thimm2019,Grant2023}, their systematic adaptation to planning domains remains an open research direction.

Important theoretical groundwork on proof systems for demonstrating unsolvability in planning has been established by~\cite{Eriksson2018}. This work proposes a rule-based proof system for proving unsolvability in planning, examining why certain planning tasks cannot be solved using \ac{STRIPS}-based approaches. While this research provides systematic approaches to proving unsolvability, it focuses on generating verifiable certificates rather than measuring degrees of inconsistency.

From a diagnostic perspective, \cite{Goebelbecker2010} introduced ``excuses'' for unsolvable planning tasks, minimal modifications required to restore solvability, such as removing specific goals or constraints. While conceptually similar to minimal inconsistent sets in logical knowledge bases~\cite{Hunter2008}, excuses focus on identifying minimal repairs rather than developing systematic inconsistency measures. Although counting excuses could provide a basic quantitative approach, this direction has not been systematically developed into an inconsistency measurement framework with comparative analysis capabilities for different types of planning conflicts.

Complementing these diagnostic approaches, methods for producing explanations that help users understand why planning tasks are unsolvable have been developed by~\cite{Sreedharan2019}, focusing on the identification of unreachable subgoals. This approach emphasizes explanatory aspects rather than quantitative measurement, but demonstrates the practical value of systematic analysis when planning problems are unsolvable. While valuable for understanding individual cases of unsolvability, these explanation methods do not provide the quantitative measures needed for systematic problem analysis, automated repair, or comparative evaluation of different unsolvable instances.

These research efforts establish important theoretical and practical foundations for understanding unsolvable planning problems. However, they reveal a fundamental limitation: existing approaches primarily focus on binary characterizations of unsolvability rather than quantitative analysis. The proof systems of~\cite{Eriksson2018} verify unsolvability, the excuses of~\cite{Goebelbecker2010} identify minimal repairs, and the explanations of~\cite{Sreedharan2019} clarify unreachable goals, but none provide the comparative quantitative measures needed for systematic analysis across different unsolvable instances. This thesis builds on these foundations by adapting inconsistency measurement techniques from propositional logic to provide systematic quantitative approaches for planning diagnostics.
